% Chapitre Introduction - Version très courte
\newpage
\chapter*{INTRODUCTION}
\addcontentsline{toc}{chapter}{Introduction}

\setcounter{page}{1}

\noindent
À l'ère du numérique, la dépendance croissante aux technologies de l'information expose nos infrastructures à des risques de sécurité considérables. L'augmentation alarmante des cyberattaques, avec des conséquences financières chiffrées en milliards d'euros, révèle l'insuffisance des systèmes de détection d'intrusions (IDS) traditionnels basés sur des signatures. Ces derniers présentent des limitations majeures, notamment leur incapacité à détecter les nouvelles formes d'attaques (zero-day).

\vspace{0.5cm}

\noindent
L'apprentissage automatique offre des perspectives prometteuses pour surmonter ces défis en permettant l'apprentissage automatique à partir des données historiques et l'adaptation aux nouveaux types de menaces. Cependant, son application à la cybersécurité soulève des défis spécifiques, principalement le déséquilibre prononcé des données où les échantillons d'attaques représentent une fraction minime du trafic total.

\vspace{0.5cm}

\noindent
Ce mémoire propose une approche innovante basée sur une stratégie hiérarchique d'apprentissage automatique pour la détection d'intrusions réseau, appliquée au dataset NSL-KDD. Notre méthodologie se distingue par une architecture en deux étapes : classification binaire (trafic légitime vs suspect) puis classification fine des types d'intrusions. Pour traiter le déséquilibre des classes, nous implémentons la technique SVM-SMOTE.

\vspace{0.5cm}

\noindent
Les contributions principales sont : (1) la démonstration de l'efficacité d'une approche hiérarchique combinant apprentissage supervisé et non-supervisé, (2) l'évaluation comparative de différents algorithmes d'apprentissage automatique, (3) la validation de l'impact des techniques de rééquilibrage sur la détection des classes minoritaires.

\vspace{0.5cm}

\noindent
Ce mémoire s'organise en quatre chapitres : méthodologie, résultats, discussion et conclusion.

\newpage